\noindent\uline{Corolario 2.9.4. Sean $P,Q,G$ grupos conmutativos que
satisfacen}
\begin{equation}\tag{2.9.4.1}
 \underline{\operatorname{Hom}}_{\mathrm{pt}}(P,G)=e,
 \qquad
 \underline{\operatorname{Hom}}_{\mathrm{pt}}(Q,G)=e
\end{equation}
\uline{(cf.\ 1.3.7). Entonces}
\sisi{\uline{se satisfacen las condiciones preliminares de 2.9.2}}%
{\uline{se satisfacen las hipótesis de la Proposición 2.9.3}},
\uline{y son válidas las conclusiones (a) y (b) de dicha proposición.
Además, si $E$ es un torsor bajo $G_{P\times Q}$ birrigidificado
respecto de $(e_P,e_Q)$, admite una estructura de biextensión de
$(P,Q)$ por $G$ compatible con sus estructuras dadas si y solo si los
dos morfismos de haces apuntados que define,}
\begin{equation}\tag{2.9.4.2}
 Q\longrightarrow R^1p_*(G_P),
 \qquad
 P\longrightarrow R^1q_*(G_Q),
\end{equation}
\uline{donde $p:P\longrightarrow e$ y $q:Q\longrightarrow e$ son
los morfismos estructurales, son morfismos de grupos.}

La primera afirmación se demostró en 1.3.7. Para la segunda, la
necesidad es clara. Para demostrar la suficiencia, se aplican
\sisi{2.9.2(b)}{la Proposición 2.9.3(b)} y 1.3.5(b). La
multiplicatividad del primer morfismo de (2.9.4.2) significa que el
torsor rigidificado $E$ bajo $G_P$ satisface la condición de 1.3.5(b)
\underline{localmente sobre $P$}. Como se observó en 1.3.7, esa
condición se cumple entonces globalmente sobre $P$, por lo que $E$
admite una estructura de extensión de $Q_P$ por $G_P$ compatible con
su estructura de torsor rigidificado. Simétricamente, $E$ admite una
estructura de extensión de $P_Q$ por $G_Q$ compatible con su
estructura de torsor rigidificado bajo $G_Q$. La conclusión se deduce
de \sisi{2.9.2(b)}{la Proposición 2.9.3(b)}.

\underline{Ejemplo 2.9.5.} Sea $S$ un esquema, sean $P$ y $Q$
esquemas abelianos sobre $S$, y sea
$G=\sisi{\underline G_{mS}}{\underline G_{m,S}}$ el esquema en grupos
multiplicativo sobre $S$. Observamos en 1.3.8 que entonces se satisface
(2.9.4.1). Un torsor birrigidificado $E$ sobre $P\times Q$ equivale a
un haz invertible birrigidificado sobre $P\times Q$. Recordemos que una
\underline{correspondencia divisorial sobre $P,Q$} es la clase de
isomorfismo de tal haz invertible birrigidificado (véase, por ejemplo,
[1]). En la topología fppf, los morfismos (2.9.4.2) son
\begin{equation}\tag{2.9.5.1}
 Q\longrightarrow\underline{\operatorname{Pic}}_{P/S},
 \qquad
 P\longrightarrow\underline{\operatorname{Pic}}_{Q/S},
\end{equation}
definidos por la correspondencia divisorial. El teorema del cuadrado
(\emph{loc.\ cit.}) dice que estos morfismos respetan necesariamente
las estructuras de grupo. Así, 2.9.4 muestra que
\underline{la categoría de biextensiones de $(P,Q)$ por
$G=\underline G_{m,S}$ es equivalente a la categoría de haces
invertibles birrigidificados sobre $P\times Q$}; es, por tanto, una
categoría rígida, y su conjunto de clases de isomorfismo se identifica
con el conjunto de correspondencias divisoriales sobre $(P,Q)$. Esta
identificación respeta manifiestamente las leyes de grupo naturales
sobre ambos conjuntos, siendo la primera la de (2.5.2.2).

\underline{Observación 2.9.6.} La aditividad del primer morfismo de
(2.9.4.2) equivale a la factorización del segundo a través del subhaz
$\underline{\operatorname{Ext}}^1(Q,G)$; obsérvese que, por 1.3.5, el
morfismo canónico
\[
 \underline{\operatorname{Ext}}^1(Q,G)
 \longrightarrow R^1q_*(G_Q)
\]
es inyectivo. Del mismo modo, la aditividad del segundo morfismo de
(2.9.4.2) significa que el primero se factoriza a través de
$\underline{\operatorname{Ext}}^1(P,G)$. Esto da isomorfismos de grupos
\begin{equation}\tag{2.9.6.1}
 \operatorname{Biext}\sisi{}{^{1}}(P,Q;G)
 \simeq
 \operatorname{Hom}\bigl(Q,\underline{\operatorname{Ext}}^1(P,G)\bigr)
 \simeq
 \operatorname{Hom}\bigl(P,\underline{\operatorname{Ext}}^1(Q,G)\bigr).
\end{equation}
En la situación de 2.9.5, estos se reducen a los isomorfismos
familiares
\begin{equation}\tag{2.9.6.2}
 \operatorname{Corr}(P,Q)
 \simeq \operatorname{Hom}(P,Q')
 \simeq \operatorname{Hom}(Q,P'),
\end{equation}
donde
\[
 P'=\underline{\operatorname{Ext}}^1(P,\underline G_m),
 \qquad
 Q'=\underline{\operatorname{Ext}}^1(Q,\underline G_m)
\]
son los esquemas abelianos duales de $P$ y $Q$, identificados,
respectivamente, con $\underline{\operatorname{Pic}}^0_{P/S}$ y
$\underline{\operatorname{Pic}}^0_{Q/S}$. En ambas fórmulas,
$\operatorname{Hom}$ denota el grupo de \underline{homomorfismos de
grupos}.
